% lecture08.tex: Lecture 8: Planetary Interiors, Structure, Composition, & Dynamics
% Planetary Systems, Kapteyn Institute

\documentclass[aspectratio=169, 11pt]{beamer}

% ── Theme and macros ────────────────────────────────────────
\usepackage{../common/beamerthemeIPS}
% macros.tex: Shared math macros for IPS lecture slides
% Mirrors definitions in book/_config.yml for consistency
% Uses \providecommand to avoid conflicts with unicode-math

% ── Derivatives ─────────────────────────────────────────────
\providecommand{\dv}[2]{\frac{\mathrm{d} #1}{\mathrm{d} #2}}
\providecommand{\dd}{}\renewcommand{\dd}{\mathrm{d}}
\providecommand{\pdv}[2]{\frac{\partial #1}{\partial #2}}

% ── Solar quantities ────────────────────────────────────────
\newcommand{\Msun}{M_{\odot}}
\newcommand{\Rsun}{R_{\odot}}
\newcommand{\Lsun}{L_{\odot}}

% ── Planetary quantities ────────────────────────────────────
\newcommand{\Mearth}{M_{\oplus}}
\newcommand{\Rearth}{R_{\oplus}}
\newcommand{\Mjup}{M_{\mathrm{J}}}
\newcommand{\Rjup}{R_{\mathrm{J}}}

% ── Constants ───────────────────────────────────────────────
\newcommand{\kB}{k_{\mathrm{B}}}

% ── Media links ─────────────────────────────────────────────
% Clickable button that opens the animated or video version of a
% figure, e.g. \mediabutton{https://...gif}{Play animation}
\providecommand{\mediabutton}[2]{\href{#1}{\beamergotobutton{#2}}}

\graphicspath{{figures/}{../common/}{../../book/08_interiors/figures/}}

% ── Metadata ────────────────────────────────────────────────
\title[Lecture 8: Planetary Interiors]{Lecture 8: Planetary Interiors\\Structure, Composition, \& Dynamics}
\subtitle{Planetary Systems}
\author[L8]{Tim Lichtenberg}
\institute{Kapteyn Astronomical Institute\\University of Groningen}
\date{2026}
\titlebgcredit{Earth's interior cross-section. Credit: NASA/JPL-Caltech}

% ════════════════════════════════════════════════════════════
\begin{document}

% ── Title slide ─────────────────────────────────────────────
\begin{frame}[plain,noframenumbering]
  \titlepage
\end{frame}

% ── Learning objectives ─────────────────────────────────────
\begin{frame}{Learning objectives}
  By the end of this lecture, you will be able to:
  \vskip8pt
  \begin{itemize}
    \item Explain how seismology, gravity fields and the moment of inertia probe interiors
    \item Derive $C/MR^2$ for uniform and two-layer planets and apply hydrostatic equilibrium with an equation of state
    \item Compare internal structures across the solar system, including the InSight results for Mars
  \end{itemize}
\end{frame}

% ── Recap ───────────────────────────────────────────────────
\begin{frame}{Recap: Lecture 7}
  \begin{itemize}
    \item Crater scaling laws and surface chronology from crater counting
    \item Volcanism, tectonics, and erosion styles across the solar system
    \item Remote sensing and cryovolcanism on outer icy worlds
  \end{itemize}
  \vskip8pt
  \textbf{Today:} What lies \textit{beneath} the surface? How do we know, and what does it tell us?
\end{frame}


% ════════════════════════════════════════════════════════════
\sectionimage[0.65]{seismic_waves}{Seismic P-wave paths through Earth. Credit: Wikimedia Commons, public domain}
\section{Probing planetary interiors}
% ════════════════════════════════════════════════════════════

% ── Three tools ──────────────────────────────────────────────
\begin{frame}{Three ways to see inside a planet}
  \begin{enumerate}\setlength\itemsep{4pt}
    \item \textbf{Seismology:} wave speeds trace density and elastic moduli
    \item \textbf{Gravity fields:} spacecraft tracking measures the mass distribution ($J_2$)
    \item \textbf{Moment of inertia:} degree of central mass concentration ($C/MR^2$)
  \end{enumerate}
  \vskip10pt
  \keyresult{For Earth all three apply; for most other worlds, gravity and $C/MR^2$ are the only probes.}
\end{frame}

% ── Seismic waves ────────────────────────────────────────────
\begin{frame}
  \ipsherokey[0.62]{seismic_waves}{P-wave paths through Earth's interior; no direct P arrival reaches the surface in the shadow zone. Credit: Wikimedia Commons, public domain.}{P-waves ($v_P = \sqrt{(K + \tfrac{4}{3}\mu)/\rho}$) travel through solids and liquids; S-waves ($v_S = \sqrt{\mu/\rho}$) need $\mu > 0$, so they \textbf{cannot propagate} through a liquid.}
\end{frame}

% ── Shadow zones hero ────────────────────────────────────────
\begin{frame}
  \ipsherokey[0.65]{prem_shadow_zones}{Rays traced through PREM, the standard seismological model of Earth's interior. Course original, from the PREM tabulation of Dziewonski \& Anderson (1981).}{No S-waves beyond 103\textdegree{} proves the outer core is liquid; P-wave refraction opens a shadow zone, 98\textdegree{} to 145\textdegree{} in this ray tracing (textbook range 103\textdegree{} to 143\textdegree{}).}
\end{frame}

% ── Gravity field ────────────────────────────────────────────
\begin{frame}{Gravity field measurements}
  \begin{itemize}\setlength\itemsep{4pt}
    \item Rotational flattening produces a quadrupole potential term $J_2$
    \item Spacecraft radio tracking detects minute Doppler velocity shifts
    \item Juno and GRAIL constrain the internal mass distribution to high precision
  \end{itemize}
\end{frame}

% ── Moment of inertia and Darwin-Radau ──────────────────────
\begin{frame}{The moment of inertia factor: $C/MR^2$}
  The moment of inertia factor measures central mass concentration:
  \vskip6pt
  \begin{columns}[c]
    \column{0.48\textwidth}
    \centering
    \footnotesize
    \begin{tabular}{l c l}
      \toprule
      \textbf{Body} & $\boldsymbol{C/MR^2}$ & \textbf{State} \\
      \midrule
      Uniform sphere & 0.400 & Homogeneous \\
      Moon           & 0.393 & Small core \\
      Mars           & 0.364 & Moderate core \\
      Mercury        & 0.346 & Large core \\
      Earth          & 0.331 & Iron core \\
      \bottomrule
    \end{tabular}
    \vskip4pt
    \sourcenote{de Pater \& Lissauer (2010)}

    \column{0.48\textwidth}
    \textbf{Darwin-Radau relation:}
    \vskip4pt\centerline{\resizebox{0.88\linewidth}{!}{$\displaystyle \frac{C}{MR^2} = \frac{2}{3}\left[1 - \frac{2}{5}\sqrt{\frac{4 - k_f}{1 + k_f}}\right]$}}\vskip4pt
    \begin{itemize}
      \item Hydrostatic bodies: from the flattening, $k_f = 3J_2/m$ with $m = \omega^2R^3/GM$
      \item Moon, Mars, Mercury: from $J_2$ and the precession or libration rate
      \item Lower value $\Rightarrow$ denser core
    \end{itemize}
  \end{columns}
\end{frame}

\begin{frame}
  \ipshero[0.76]{moment_of_inertia_factors}{The moment of inertia factor $C/MR^2$ from the uniform sphere (0.400) down to Earth (0.331), in the order of central mass concentration. Values after de Pater \& Lissauer (2010). Course figure.}
\end{frame}

% ── MoI factor hero ──────────────────────────────────────────
\begin{frame}
  \ipshero{moi_factor}{Moment of inertia factors for nine solar-system bodies, all below the uniform-sphere value of 0.400. Course original; lunar value from Williams et al.\ (2014).}
\end{frame}


% ════════════════════════════════════════════════════════════
\sectionimage{terrestrial_interiors_comparison}{Differentiated interiors of the terrestrial bodies. Credit: NASA Solar System Exploration, public domain}
\section{Blackboard derivation: moment of inertia factor}
% ════════════════════════════════════════════════════════════

% ── Setup ────────────────────────────────────────────────────
\begin{frame}{Deriving $C/MR^2$}
  \begin{columns}[c]
    \column{\recapfigcolnarrow}
    \ipsrecapfig{moi_shell_sketch}
    \source{Course original}

    \column{\recaptextcolwide}
    \textbf{Goal:} Derive the moment of inertia factor for (a) a uniform sphere and (b) a two-layer differentiated body.
    \vskip6pt
    \textbf{Starting point:} for a spherically symmetric body,
    $$C = \frac{8\pi}{3}\int_0^R \rho(r)\,r^4\,\dd r$$
    A mass element at colatitude $\theta$ contributes $r_\perp^2\,\dd m = r^2\sin^2\theta\,\dd m$; over a thin shell, $\sin^2\theta \to 2/3$.
  \end{columns}
\end{frame}

% ── Uniform sphere and two-layer formulas ───────────────────
\begin{frame}{Summary of derived results}
  \textbf{Uniform sphere:} $\rho = \text{const}$
  $$C = \frac{8\pi}{15}\rho R^5, \quad M = \frac{4}{3}\pi\rho R^3 \quad \implies \quad \frac{C}{MR^2} = \frac{2}{5} = 0.400$$
  \vskip6pt
  \textbf{Two-layer body:} core ($\rho_c, R_c$) and mantle ($\rho_m, R$), with $x \equiv R_c/R$ and $f \equiv \rho_c/\rho_m$:
  \keyresult{$$\frac{C}{MR^2} = \frac{2}{5}\cdot\frac{(f-1)\,x^5 + 1}{(f-1)\,x^3 + 1}$$}
  \vskip4pt
  Since $x^5 < x^3$ for $0 < x < 1$, any dense core ($f > 1$) pulls $C/MR^2 < 0.400$.
\end{frame}

% ── Two-layer bodies visualised ──────────────────────────────
\begin{frame}
  \ipsherokey[0.65]{moi_two_layer}{Two-layer body at fixed $f = \rho_c/\rho_m = 2.5$: $C/MR^2$ against the core radius fraction $x = R_c/R$. Course original.}{$C$ weights mass by $r^2$, so mass moved inward lowers $C$ but not $M$: Earth ($x = 0.546$) gives $C/MR^2 \approx 0.345$, and the gap to the measured 0.331 is self-compression.}
\end{frame}


% ════════════════════════════════════════════════════════════
\sectionimage[0.65]{jupiter_interior}{Jupiter's interior. Credit: NASA/JPL-Caltech/SwRI/John E. Connerney (PIA25062)}
\section{Equations of state}
% ════════════════════════════════════════════════════════════

% ── Hydrostatic equilibrium ──────────────────────────────────
\begin{frame}{Hydrostatic equilibrium in spherical symmetry}
  Pressure at each depth supports the weight of the overlying material:
  \vskip8pt
  \keyresult{$$\dv{P}{r} = -\rho(r)\,g(r) = -\frac{G\,m(r)\,\rho(r)}{r^2}$$}
  \vskip6pt
  where $m(r) = 4\pi\int_0^r \rho(r')\,r'^2\,\dd r'$ is the mass enclosed within radius $r$.
  \vskip8pt
  With an equation of state $P = P(\rho, T)$, this integrates inward to the pressure-density profile.
\end{frame}

% ── Central pressure ─────────────────────────────────────────
\begin{frame}
  \ipsherokey[0.6]{central_pressure}{Central pressure from interior models against the uniform-density estimate for six solar-system bodies. Course original.}{Uniform density gives $P_c \approx 3GM^2/(8\pi R^4)$, a lower bound: about 170~GPa against the actual 360~GPa for Earth, 25 against 40~GPa for Mars; real planets sit above the line because density increases inward.}
\end{frame}

% ── PALEOS EOS hero ──────────────────────────────────────────
\begin{frame}
  \ipshero{eos_paleos}{Equations of state for iron, silicate and water on a cold ($\sim$300~K) and a hot ($\sim$4000~K) isotherm; the density jumps mark phase transitions. Course original, from the PALEOS tables of Attia et al.\ (2026).}
\end{frame}

% ── Birch's law ──────────────────────────────────────────────
\begin{frame}
  \ipsherokey[0.61]{birchs_law}{Birch's law: $v_P$ against density for mantle minerals, trend after Birch (1952). Course original.}{$v_P = a(\bar{M}) + b\,\rho$: at a given mean atomic weight $\bar{M}$, compressional velocity scales linearly with density, so seismic speeds map to density. Iron-rich minerals plot below the trend, so the law needs a known composition.}
\end{frame}

% ── Adams-Williamson equation ────────────────────────────────
\begin{frame}{Adams-Williamson equation}
  Seismic wave speeds integrate to a density profile under adiabatic self-compression:
  \vskip4pt
  $$\dv{\rho}{r} = -\frac{\rho(r)\,g(r)}{\phi(r)}, \quad \phi = v_P^2 - \frac{4}{3}v_S^2$$
  \vskip4pt
  \begin{itemize}\setlength\itemsep{3pt}
    \item Seismic parameter $\phi \equiv K_S/\rho = v_P^2 - \tfrac{4}{3}v_S^2$
    \item Assumes hydrostatic equilibrium, an adiabatic gradient and uniform composition
    \item Discontinuities mark phase changes or compositional boundaries (e.g., the core-mantle boundary)
  \end{itemize}
\end{frame}

% ── Adams-Williamson hero ────────────────────────────────────
\begin{frame}
  \ipshero[0.76]{adams_williamson}{The Adams-Williamson relation against PREM: adiabatic self-compression matches the lower mantle, and the density jump at the core-mantle boundary exposes a change in composition. Course original, from the PREM tabulation of Dziewonski \& Anderson (1981).}
\end{frame}


% ════════════════════════════════════════════════════════════
\sectionimage[0.65]{earth_interior}{Earth's interior. Credit: NASA/JPL-Caltech/SwRI/John E. Connerney (PIA25063)}
\section{Earth's interior}
% ════════════════════════════════════════════════════════════

% ── Earth cross-section ──────────────────────────────────────
\begin{frame}
  \ipsherokey[0.62]{earth_interior}{Earth's interior in cross-section. Credit: NASA/JPL-Caltech.}{\textbf{Crust} (0--70~km), \textbf{mantle} (Moho to 2891~km), liquid \textbf{outer core} (2891--5150~km) and solid \textbf{inner core} (5150--6371~km, where the melting point at 330--360~GPa exceeds 5000--6000~K); mean density 5515~kg\,m$^{-3}$, $C/MR^2 = 0.3307$, central pressure about 360~GPa.}
\end{frame}

% ── PREM ─────────────────────────────────────────────────────
\begin{frame}
  \ipsherokey[0.61]{prem_density}{Density against radius in the \textbf{Preliminary Reference Earth Model} (PREM). Course original, evaluated from the coefficients of Dziewonski \& Anderson (1981).}{Density rises from $\sim$3300~kg\,m$^{-3}$ in the upper mantle to $\sim$5570 at its base, jumps to 9900--12{,}200 in the outer core and reaches $\sim$13{,}000 in the inner core. The 410 and 660~km steps are phase transitions; the core-mantle boundary is silicate to iron.}
\end{frame}

% ── PREM full profile hero ───────────────────────────────────
\begin{frame}
  \ipshero{dziewonski_anderson1981_prem_model}{PREM radial profiles of wave speeds, density and anisotropy; $v_S$ vanishes in the liquid outer core and resumes in the solid inner core. Credit: Dziewonski \& Anderson (1981), Fig.~8.}
\end{frame}

% ── Crust, mantle, and core ──────────────────────────────────
\begin{frame}{Crust, mantle, and core}
  \centering
  \small
  \begin{tabular}{l c c l}
    \toprule
    \textbf{Layer} & \textbf{Depth (km)} & $\boldsymbol{\rho}$ \textbf{(kg\,m$^{-3}$)} & \textbf{Composition} \\
    \midrule
    Crust           & 0--70      & 2700--3000     & Silicates (Moho base) \\
    Upper mantle    & 70--410    & 3300--3700     & Olivine, pyroxene \\
    Transition zone & 410--660   & 3700--4400     & Wadsleyite, ringwoodite \\
    Lower mantle    & 660--2891  & 4400--5500     & Bridgmanite \\
    Outer core      & 2891--5150 & 9900--12{,}200 & Liquid Fe-Ni \\
    Inner core      & 5150--6371 & $\sim$13{,}000  & Solid Fe-Ni \\
    \bottomrule
  \end{tabular}
  \vskip8pt
  \keyresult{Convection in the liquid outer core powers the geodynamo, fed by inner-core growth: crystallising iron releases latent heat and light elements.}
\end{frame}


% ════════════════════════════════════════════════════════════
\sectionimage{earth_heat_flow_map}{Global surface heat flow, the surface expression of mantle convection. Credit: Fabio Crameri, Wikimedia Commons, CC BY-SA 4.0}
\section{Mantle rheology \& convection}
% ════════════════════════════════════════════════════════════

% ── Solid-state creep ────────────────────────────────────────
\begin{frame}
  \ipsherokey[0.66]{creep_viscosity_ranges}{Mantle viscosity by creep mechanism. Course schematic.}{Solid rock flows once $T > 0.5\,T_{\mathrm{melt}}$: \textbf{dislocation creep} (defect glide) in the upper mantle at $10^{20}$ to $10^{21}$~Pa\,s, \textbf{diffusion creep} (vacancy migration) in the lower mantle at $10^{22}$ to $10^{23}$~Pa\,s.}
\end{frame}

% ── Rayleigh number ──────────────────────────────────────────
\begin{frame}{The Rayleigh number}
  Thermal buoyancy against viscous resistance decides whether a mantle convects:
  \vskip4pt
  $$\mathrm{Ra} = \frac{\alpha\,\rho\,g\,\Delta T\,d^3}{\kappa\,\eta}$$
  \vskip4pt
  \begin{itemize}\setlength\itemsep{3pt}
    \item $\mathrm{Ra}$ is the ratio of buoyancy driving forces to viscous dissipation
    \item Critical threshold for the onset of convection: $\mathrm{Ra}_c \sim 10^3$
    \item Earth's mantle: $\mathrm{Ra} \sim 10^7\text{--}10^8 \gg \mathrm{Ra}_c$, vigorous convection
  \end{itemize}
\end{frame}

\begin{frame}
  \ipshero[0.76]{rayleigh_viscosity}{The Rayleigh number against mantle viscosity for the course's Earth-mantle values: Earth at $\eta = 10^{21}$ to $10^{22}$~Pa\,s sits at $\mathrm{Ra} \approx 10^7$ to $10^8$, four to five orders of magnitude above $\mathrm{Ra}_c \approx 10^3$. Course figure.}
\end{frame}

% ── Convection regimes hero ──────────────────────────────────
\begin{frame}
  \ipshero[0.76]{convection_regimes}{End-member regimes of mantle convection; Earth convects mainly as a whole mantle, with the 660~km discontinuity a partial, regionally variable barrier. Schematic; boundary radius exaggerated.}
\end{frame}


% ════════════════════════════════════════════════════════════
\sectionimage[0.65]{earth_interior}{Earth's interior. Credit: NASA/JPL-Caltech/SwRI/John E. Connerney (PIA25063)}
\section{Phase transitions in the mantle}
% ════════════════════════════════════════════════════════════

% ── Phase transition table ───────────────────────────────────
\begin{frame}{Mineral phase transitions with depth}
  \centering
  \small
  \begin{tabular}{c l c}
    \toprule
    \textbf{Depth (km)} & \textbf{Transition} & \textbf{$\Delta\rho$} \\
    \midrule
    410  & Olivine $\rightarrow$ wadsleyite & $\sim$5\% \\
    520  & Wadsleyite $\rightarrow$ ringwoodite & $\sim$2\% \\
    660  & Ringwoodite $\rightarrow$ bridgmanite + ferropericlase & $\sim$8\% \\
    $\sim$2700 & Bridgmanite $\rightarrow$ post-perovskite & $\sim$1--2\% \\
    \bottomrule
  \end{tabular}
  \vskip10pt
  \raggedright
  \begin{itemize}\setlength\itemsep{3pt}
    \item At 410 and 520~km olivine ($\mathrm{(Mg,Fe)_2SiO_4}$) turns into denser polymorphs of the same composition; at 660~km ringwoodite splits into bridgmanite and ferropericlase
    \item Driven by pressure: atoms rearrange into denser packing
    \item These transitions produce the primary seismic discontinuities in PREM
  \end{itemize}
\end{frame}

% ── Mantle phase diagram hero ────────────────────────────────
\begin{frame}
  \ipshero[0.77]{mantle_phase_diagram}{Phase boundaries of the mantle against an approximate geotherm; seismic discontinuities occur where the geotherm crosses a phase boundary. After Turcotte \& Schubert (2002) and Murakami et al.\ (2004).}
\end{frame}

% ── 660 km discontinuity ─────────────────────────────────────
\begin{frame}{The 660~km discontinuity}
  \begin{itemize}\setlength\itemsep{4pt}
    \item Ringwoodite decomposes into bridgmanite plus ferropericlase with an $\sim$8\% density jump
    \item Negative Clapeyron slope ($\dd P/\dd T < 0$): in a cold slab the boundary sits deeper, so buoyant ringwoodite persists below 660~km
    \item The upward buoyancy is a partial barrier to whole-mantle convection
  \end{itemize}
  \vskip8pt
  \keyresult{The endothermic ringwoodite breakdown at 660~km resists slab penetration.}
\end{frame}


% ── Break ───────────────────────────────────────────────────
\breakslide


% ════════════════════════════════════════════════════════════
\sectionimage[0.65]{mercury_interior}{Mercury's layered interior. After Margot et al.\ (2018)}
\section{Comparative terrestrial interiors}
% ════════════════════════════════════════════════════════════

% ── The Moon ─────────────────────────────────────────────────
\begin{frame}
  \ipsherokey[0.6]{weber2011_lunar_core}{The lunar interior from reanalysed Apollo seismograms: a small metallic core with a liquid outer and a solid inner part. Credit: Weber et al.\ (2011), Fig.~1A.}{$C/MR^2 = 0.393$ and $\bar{\rho} = 3346$~kg\,m$^{-3}$, close to uncompressed mantle rock: the core is only 200--380~km in radius ($R_c/R \sim 0.1$--0.2, $<2$\% of the mass), as expected for a Moon formed by a giant impact from Earth's mantle.}
\end{frame}

% ── Mars interior hero ───────────────────────────────────────
\begin{frame}
  \ipsherokey[0.61]{mars_interior_insight}{Mars's interior from InSight seismology: a 24--72~km crust, a mantle extending to $\sim$1560~km depth, and the core. Credit: NASA/JPL-Caltech (PIA25282).}{The seismic boundary at $R_c \sim 1830$~km was reinterpreted in 2023 as the top of a molten silicate layer above a smaller ($\sim$1650--1675~km) liquid iron core of about 6500~kg\,m$^{-3}$.}
\end{frame}

% ── Mercury ──────────────────────────────────────────────────
\begin{frame}
  \ipsherokey[0.61]{mercury_interior}{Mercury's layered interior from MESSENGER gravity and libration data. After Margot et al.\ (2018).}{$C/MR^2 = 0.346$ and $\bar{\rho} = 5427$~kg\,m$^{-3}$, the highest uncompressed density of the terrestrial planets: the metallic core reaches $R_c \sim 2000$--2050~km ($\sim$83\% of the radius, $\sim$74\% of the mass), and the weak dynamo indicates a liquid outer core.}
\end{frame}

% ── Comparative table ────────────────────────────────────────
\begin{frame}{Comparing terrestrial body interiors}
  \centering
  \small
  \begin{tabular}{l c c c c}
    \toprule
    \textbf{Body} & $\boldsymbol{\bar{\rho}}$ \textbf{(kg\,m$^{-3}$)} & $\boldsymbol{C/MR^2}$ & $\boldsymbol{R_c/R}$ & \textbf{Dynamo?} \\
    \midrule
    Moon    & 3346 & 0.393 & $\sim$0.1--0.2 & No (ancient) \\
    Mars    & 3934 & 0.364 & $\sim$0.49--0.54 & No (ancient) \\
    Venus   & 5243 & $\sim$0.33 & $\sim$0.55 & No \\
    Earth   & 5515 & 0.331 & 0.55 & Yes \\
    Mercury & 5427 & 0.346 & $\sim$0.83 & Weak \\
    \bottomrule
  \end{tabular}
  \vskip4pt
  {\tiny\color{ipsText!60} Mars: upper value from Stähler et al.\ (2021), lower from the 2023 reanalyses (Khan et al.; Samuel et al.); Mercury: Margot et al.\ (2018)}
  \vskip6pt
  \raggedright
  \begin{itemize}\setlength\itemsep{3pt}
    \item Higher \emph{uncompressed} density $\Rightarrow$ larger iron core fraction; self-compression lifts Earth's mean density above Mercury's
    \item Lower $C/MR^2$ $\Rightarrow$ stronger central mass concentration
    \item A dynamo requires a convecting liquid conductor
  \end{itemize}
\end{frame}

% ── Terrestrial comparison hero ──────────────────────────────
\begin{frame}
  \ipshero{terrestrial_interiors_comparison}{The terrestrial bodies in cross-section at relative scale: a silicate mantle over a metallic core, with core fractions from $\sim$10--20\% to $\sim$83\%. Credit: NASA Solar System Exploration, public domain.}
\end{frame}


% ════════════════════════════════════════════════════════════
\sectionimage{saturn_interior}{Saturn's interior. Credit: NASA/JPL-Caltech}
\section{Giant planet interiors}
% ════════════════════════════════════════════════════════════

% ── Jupiter interior hero ────────────────────────────────────
\begin{frame}
  \ipshero{jupiter_interior}{Jupiter's interior: molecular and metallic hydrogen above an extended, dilute core enriched in heavy elements. Credit: NASA/JPL-Caltech/SwRI/John E. Connerney (PIA25062).}
\end{frame}

\begin{frame}
  \ipshero[0.76]{wahl2017_jupiter_density}{Density profiles of Jupiter interior models from Juno gravity: dilute-core solutions spread heavy elements out to half the radius, against classical compact cores. Credit: Wahl et al.\ (2017), Fig.~1.}
\end{frame}

% ── Saturn: helium rain ──────────────────────────────────────
\begin{frame}
  \ipsherokey[0.61]{saturn_interior}{Saturn's interior, with a diffuse heavy-element core that ring seismology extends to $\sim$60\% of the radius. Credit: NASA/JPL-Caltech.}{As Saturn cools, helium becomes immiscible in metallic hydrogen at $P \sim 2$~Mbar and rains toward the centre; the released gravitational energy, with Kelvin-Helmholtz contraction, explains why Saturn radiates $\sim$1.8$\times$ the solar flux it absorbs.}
\end{frame}

% ── Helium rain hero ─────────────────────────────────────────
\begin{frame}
  \ipshero[0.76]{mankovich_fortney2020_jupiter_saturn_schematic}{Helium phase separation in Jupiter and Saturn; sinking helium droplets release gravitational potential energy and explain Saturn's excess luminosity. Credit: Mankovich \& Fortney (2020), Fig.~3.}
\end{frame}


% ════════════════════════════════════════════════════════════
\sectionimage{uranus_clouds_voyager}{Uranus from Voyager 2. Credit: NASA/JPL-Caltech}
\section{Ice giant interiors}
% ════════════════════════════════════════════════════════════

% ── Ice giant structures hero ────────────────────────────────
\begin{frame}
  \ipsherokey{helled2020_icegiant_structures}{Four candidate internal structures of an ice giant. Credit: Helled et al.\ (2020), Fig.~4.}{Canonical layers: an H/He envelope ($\sim$20\% of the mass), a fluid ice mantle ($\sim$60--70\%: $\mathrm{H_2O}$, $\mathrm{NH_3}$, $\mathrm{CH_4}$ at 10--500~GPa, possibly superionic) and a rocky core; gravity alone cannot separate sharp layers from gradients.}
\end{frame}

% ── Unusual magnetic fields ──────────────────────────────────
\begin{frame}
  \ipsherokey[0.61]{soderlund2020_planetary_field_morphologies}{Radial magnetic field at the surface of six planets in Mollweide projection. Credit: Soderlund et al.\ (2020), Fig.~1.}{Uranus and Neptune generate multipolar fields tilted by 59\textdegree{} and 47\textdegree{} and offset from the centre, from dynamos in thin outer fluid ice shells, unlike the axial dipoles of Earth, Jupiter and Saturn.}
\end{frame}


% ════════════════════════════════════════════════════════════
\sectionimage{europa_interior}{Europa's subsurface ocean. Credit: NASA/JPL-Caltech (PIA26438)}
\section{Icy moon interiors \& ocean worlds}
% ════════════════════════════════════════════════════════════

% ── Ocean worlds overview ────────────────────────────────────
\begin{frame}{Ocean worlds: liquid water beneath ice}
  Tidal flexing and radiogenic decay sustain subsurface oceans beneath insulating ice shells:
  \vskip6pt
  \centering
  \small
  \begin{tabular}{l c c l}
    \toprule
    \textbf{Moon} & $\boldsymbol{R}$ \textbf{(km)} & $\boldsymbol{\rho}$ \textbf{(kg\,m$^{-3}$)} & \textbf{Ocean evidence} \\
    \midrule
    Europa    & 1561 & 3013 & Magnetic induction, chaos terrain \\
    Enceladus & 252  & 1609 & Plumes, libration \\
    Titan     & 2575 & 1882 & Tidal Love number $k_2$ \\
    Ganymede  & 2634 & 1936 & Auroral oscillations \\
    \bottomrule
  \end{tabular}
\end{frame}

% ── Europa interior hero ─────────────────────────────────────
\begin{frame}
  \ipshero{europa_interior}{Europa's subsurface environment: an ice shell over a global ocean roughly 100~km deep in contact with a rocky seafloor. Credit: NASA/JPL-Caltech (PIA26438).}
\end{frame}

% ── Tidal resonance hero ─────────────────────────────────────
\begin{frame}
  \ipshero{tidal_resonance}{The Io-Europa-Ganymede 4:2:1 Laplace resonance forces the eccentricities whose tidal flexing keeps Europa's ocean liquid. Course original; orbital periods from NASA/JPL Solar System Dynamics.}
\end{frame}

% ── Enceladus interior hero ──────────────────────────────────
\begin{frame}
  \ipshero[0.77]{enceladus_interior}{Enceladus in cutaway: a global ocean between the ice shell and a porous core feeds hydrothermal venting and south-polar plumes radiating $\sim$16~GW. Credit: NASA/JPL-Caltech (PIA19656).}
\end{frame}

% ── Titan interior hero ──────────────────────────────────────
\begin{frame}
  \ipshero[0.76]{titan_interior}{Titan's interior from Cassini gravity and radio science: a large tidal Love number $k_2$ indicates an ocean that decouples the outer ice shell from the high-pressure ice and rock below. Credit: NASA/JPL.}
\end{frame}


% ════════════════════════════════════════════════════════════
\sectionimage{mars_interior_insight}{Mars's interior from InSight. Credit: NASA/JPL-Caltech (PIA25282)}
\section{Recent advances}
% ════════════════════════════════════════════════════════════

% ── Stähler wedge hero ───────────────────────────────────────
\begin{frame}
  \ipshero[0.76]{stahler2021_mars_wedge}{InSight seismic ray paths across Mars delineate the mantle and the core shadow zone; the 2023 reanalyses put a molten silicate layer above a $\sim$1650--1675~km core. Credit: St\"ahler et al.\ (2021), Fig.~3.}
\end{frame}

% ── Mars velocity profiles hero ──────────────────────────────
\begin{frame}
  \ipshero[0.76]{stahler2021_mars_vp_vs}{Mars velocity profiles from InSight: the mantle $v_S$ and $v_P$ profiles end at the core boundary near 1560~km depth, below which no shear wave propagates, a liquid metallic core. Credit: St\"ahler et al.\ (2021), Fig.~2A.}
\end{frame}

% ── Summary ────────────────────────────────────────────────
\begin{frame}{Summary}
  \begin{enumerate}\setlength\itemsep{6pt}
    \item Seismology, gravity fields and $C/MR^2$ probe deep interiors: 0.400 for a uniform sphere, lower values mean differentiation
    \item Hydrostatic equilibrium and equations of state set the radial profiles; mantle phase transitions at 410 and 660~km shape convection
    \item Tidal heating maintains subsurface oceans across the icy worlds
  \end{enumerate}
  \vskip8pt
  \textbf{Next: Lecture 9} turns to Earth and Venus: twins in size, contrasts in evolution.
\end{frame}

% ── Final slide ─────────────────────────────────────────────
\begin{frame}[plain,noframenumbering]
  \begin{tikzpicture}[remember picture, overlay]
    \ipscontain{title_background}
    \fill[ipsPrimary, opacity=0.70] (current page.north west) rectangle (current page.south east);
  \end{tikzpicture}
  \vfill
  \begin{center}
    {\color{white}\Large\bfseries Questions?}
    \vskip20pt
    {\color{white!85}\large Next: Lecture 9, Rocky planets I: Earth \& Venus}\\[4pt]
    {\color{white!85}\small Mon 5 Oct, 15:00--17:00, room 5419.0161 (Collegezaal)}\\[12pt]
    {\color{white!85}\small Next tutorial: Fri 2 Oct, 13:00--15:00, room 5419.0161 (Collegezaal)}\\[20pt]
    {\color{white!70}\small \texttt{https://ips.formingworlds.space}}
  \end{center}
  \vfill
\end{frame}

\end{document}
